% =====================================================================
%  Report - Project 2: FIR filter on FPGA (Red Pitaya)
%  Build (from the rapport/ folder):
%     pdflatex rapport.tex && pdflatex rapport.tex
%  (2 passes for the table of contents and references)
% =====================================================================
\documentclass[11pt,a4paper]{article}

% ---------- Language / encoding ----------
\usepackage[T1]{fontenc}
\usepackage[utf8]{inputenc}
\usepackage{lmodern}
\usepackage[english]{babel}

% ---------- Layout ----------
\usepackage[margin=2.2cm]{geometry}
\usepackage{fancyhdr}
\usepackage{graphicx}
\usepackage{float}
\usepackage{subcaption}
\usepackage{booktabs}
\usepackage{amsmath}
\usepackage[per-mode=symbol]{siunitx}
\sisetup{retain-explicit-plus}
\usepackage{xcolor}
\usepackage{tikz}
\usetikzlibrary{arrows.meta,positioning}
\usepackage{listings}
\usepackage[hidelinks]{hyperref}

\graphicspath{{figures/}}
\setlength{\emergencystretch}{3em}
% Allow line breaks after underscores in long signal names
\let\oldunderscore\_
\renewcommand{\_}{\oldunderscore\allowbreak}

% ---------- Header / footer ----------
\setlength{\headheight}{14pt}
\pagestyle{fancy}
\fancyhf{}
\lhead{FPGA -- Project 2: FIR}
\rhead{ELISE M2}
\cfoot{\thepage}

% ---------- VHDL code style ----------
\definecolor{kw}{RGB}{0,0,160}
\definecolor{cm}{RGB}{0,120,0}
\definecolor{bg}{RGB}{247,247,247}
\lstdefinestyle{vhdl}{
  language=VHDL,
  basicstyle=\ttfamily\footnotesize,
  keywordstyle=\color{kw}\bfseries,
  commentstyle=\color{cm}\itshape,
  backgroundcolor=\color{bg},
  numbers=left, numberstyle=\tiny\color{gray},
  frame=single, breaklines=true, showstringspaces=false,
  tabsize=4, columns=fullflexible, keepspaces=true
}
\lstdefinestyle{sh}{
  basicstyle=\ttfamily\footnotesize, backgroundcolor=\color{bg},
  frame=single, breaklines=true, columns=fullflexible
}
% Allows accented characters in code comments
\lstset{literate={é}{{\'e}}1 {è}{{\`e}}1 {à}{{\`a}}1 {ç}{{\c c}}1 {ê}{{\^e}}1 {ù}{{\`u}}1}

% Inserts a figure if the file exists, otherwise a "to add" placeholder box
% usage: \fig{file_name}{caption}{label}
\newcommand{\fig}[3]{%
  \begin{figure}[H]\centering
    \IfFileExists{figures/#1}
      {\includegraphics[width=0.85\linewidth]{#1}}
      {\fbox{\parbox[c][4cm][c]{0.8\linewidth}{\centering\color{gray}
        Figure to add: \texttt{figures/\detokenize{#1}}}}}
    \caption{#2}\label{#3}
  \end{figure}}

% Marker for text still to write (delete once done)
\newcommand{\todo}[1]{\textcolor{red}{[To write: #1]}}


% =====================================================================
\begin{document}

% ---------------------------- Title page ---------------------------
\begin{titlepage}
  \centering
  {\large Université Marie et Louis Pasteur\par}
  {\large ELISE Master -- 2\textsuperscript{nd} year\par}
  \vspace{3cm}
  {\Huge\bfseries Project 2: FIR\par}
  \vspace{0.5cm}
  {\Large Finite Impulse Response filters on FPGA\\(Red Pitaya, Zynq-7010)\par}
  \vspace{3cm}
  {\large \textbf{Authors:} LAST NAME First name -- LAST NAME First name\par}
  \vspace{0.5cm}
  {\large \textbf{Teacher:} NAME\par}
  \vfill
  {\large Academic year 2026--2027\par}
\end{titlepage}

\tableofcontents
\newpage

% =====================================================================
\section{Introduction}
\todo{goal of the project: FIR filter at $f_s = 125/8 = \SI{15.625}{\mega\hertz}$,
low-pass in VHDL, band-pass with the FIR Compiler IP, PS/PL switching of the two
filters from Linux. Recall the FIR equation $y(n) = \sum_k b_k\,x(n-k)$.}

% =====================================================================
\section{Low-pass filter coefficients (Octave)}
\todo{fir1 call: order 4, $f_c = \SI{0.5}{\mega\hertz}$, $f_s = 125/8$ MHz;
script \texttt{filter.m}; coefficients obtained}

\fig{bode_lowpass.png}{Frequency response of the 4th-order low-pass filter (freqz)}{fig:bodelp}

\paragraph{Real cutoff frequency.}
\todo{measured $-6$ dB frequency and why it is not \SI{0.5}{\mega\hertz}
(only 5 coefficients, wide transition band)}

% =====================================================================
\section{VHDL description of the low-pass filter}

\subsection{Integer coefficients on 8 bits}
\todo{lecture method: $n_F = 8 + 1 = 9$, scale $2^9 = 512$; table float /
$\times 512$ / integer; sum = 512 (gain 1)}

\begin{table}[H]
\centering
\begin{tabular}{cccccc}
\toprule
 & $b_0$ & $b_1$ & $b_2$ & $b_3$ & $b_4$ \\
\midrule
Octave            & & & & & \\
$\times 512$      & & & & & \\
Integer           & & & & & \\
\bottomrule
\end{tabular}
\caption{Quantization of the coefficients}\label{tab:coef}
\end{table}

\subsection{Filter structure}
\todo{entity (generics b0..b4, ports clk, enable, x, y on 14 bits unsigned);
delay line x0..x4, products p0..p4 (22 bits), sum (23 bits), output
\texttt{sum(22 downto 9)}; synchronous enable as in the lecture; latency}

\fig{fir_structure.png}{Structure of the FIR filter}{fig:firstruct}

\subsection{Testbench and simulation}
\todo{testbench: clock 64 ns, constant inputs first (impulse 512, constant 16383),
then sine inputs at 100 kHz, 500 kHz, 3.31 MHz, 5 MHz}

\fig{sim_impulse_step.png}{Simulation: impulse (512) and constant input (16383)}{fig:simimp}
\fig{sim_sines.png}{Simulation with sine inputs (same vertical scale for $x$ and $y$)}{fig:simsin}

\begin{table}[H]
\centering
\begin{tabular}{ccccc}
\toprule
Frequency & $x$ (peak to peak) & $y$ (peak to peak) & Simulated gain & Octave gain \\
\midrule
\SI{100}{\kilo\hertz} & & & & \\
\SI{500}{\kilo\hertz} & & & & \\
\SI{3.31}{\mega\hertz} & & & & \\
\SI{5}{\mega\hertz} & & & & \\
\bottomrule
\end{tabular}
\caption{Simulated attenuation compared with Octave}\label{tab:simgain}
\end{table}

\paragraph{Why 14 bits are enough at the output.}
\todo{sum of coefficients = 512 = $2^9$, gain $\leq 1$, max sum
$16383 \times 512 < 2^{23}$, shift by 9 bits}

% =====================================================================
\section{Encapsulation and test on the board}

\subsection{Encapsulating file}
\todo{based on ADC\_DAC.vhd: 8-state counter at 125 MHz, enable at 125/8 MHz,
ADC channel A to the filter, filter to the DAC}

\subsection{Synthesis}
\todo{openXC7 commands; SITEWIRE/DSP48 error if it happens}

\subsection{First test: shifted sine}
\todo{generator 0.5 Vpp ($-0.25$ to $+0.25$ V); observation: truncated sine with
shifted parts; explanation: +1 V gives code 0 and $-1$ V gives $2^{14}-1$}

\fig{board_shifted.png}{Output before correcting the ADC code}{fig:shifted}

\subsection{Correction of the ADC/DAC coding}
\todo{correction in the encapsulating file and why it works}

\subsection{Low-pass behaviour on the board}
\todo{measurements at several frequencies, compare with Octave}

\fig{board_lowpass.png}{Low-pass filter on the board}{fig:boardlp}

\begin{table}[H]
\centering
\begin{tabular}{cccc}
\toprule
Frequency & Input (Vpp) & Output (Vpp) & Gain (dB) \\
\midrule
 & & & \\
\bottomrule
\end{tabular}
\caption{Measured attenuation of the low-pass filter}\label{tab:boardlp}
\end{table}

% =====================================================================
\section{Band-pass filter with the FIR Compiler IP}

\subsection{Coefficients (Octave)}
The script \texttt{filter\_pb.m} computes the 31 coefficients with
\texttt{b = fir1(30, [2 2.5]/(fs/2), 'bandpass')}, $f_s = 125/8 = \SI{15.625}{\mega\hertz}$.
The coefficients are symmetric (linear phase) and their sum is close to 0,
so the gain at DC is almost zero (sum of the integer coefficients: 7, i.e. $7/512 \approx 0.014$).

\paragraph{Real cutoff frequencies.}
With \texttt{freqz} (figure~\ref{fig:bodebp}), the $-6$~dB band goes from
\SI{1.750}{\mega\hertz} to \SI{2.734}{\mega\hertz} instead of 2 to
\SI{2.5}{\mega\hertz}. As for the low-pass filter, \texttt{fir1} uses a window
method, so the cutoff frequencies are not placed exactly at $-6$~dB. The
requested band is very narrow (\SI{0.5}{\mega\hertz}, about 3\,\% of $f_s$),
and 31 coefficients are not enough to get sharp edges: the transition bands
are wide and the pass band becomes larger than requested.

\fig{bode_bandpass.png}{Frequency response of the 30th-order band-pass filter}{fig:bodebp}

\subsection{FIR Compiler configuration}
The FIR Compiler IP is configured with:
\begin{itemize}
  \item filter options: coefficient vector filled with the 31 floating-point
        values from Octave;
  \item channel specification: input sampling frequency and clock frequency
        at \SI{15.625}{\mega\hertz} (one sample per clock cycle);
  \item implementation: signed coefficients on 8 bits with ``best precision
        fraction length'', signed input data on 14 bits, full precision
        output.
\end{itemize}
The input \texttt{s\_axis\_data\_tvalid} is set permanently to `1' with a
Constant IP, so that the filter takes a new sample at each clock edge.
With ``best precision'', Vivado chooses 9 fractional bits for the
coefficients: the integer coefficients are $\mathrm{round}(b \times 2^9)$,
which is the same scaling as for the low-pass filter:
\[
2,\ 4,\ 4,\ -2,\ -11,\ -16,\ -6,\ 18,\ 36,\ 28,\ -10,\ -49,\ -54,\ -15,\ 41,\
67,\ 41,\ \dots,\ 4,\ 4,\ 2
\]
The full precision output is on 24 bits and the latency given in the summary
tab is 22 clock cycles (about \SI{1.4}{\micro\second}).

\fig{fir_compiler.png}{Block design of the band-pass filter: FIR Compiler with \texttt{s\_axis\_data\_tvalid} set to `1' by a Constant IP}{fig:fircomp}

\subsection{Simulation in Vivado}
The testbench \texttt{tb\_fir.vhd} instantiates the block design with a
\SI{64}{\nano\second} clock (\SI{15.625}{\mega\hertz}). The input is a
14-bit signed value, sign-extended on the 16 bits of
\texttt{s\_axis\_data\_tdata}. The test has two parts:
\begin{itemize}
  \item an impulse of 8191 during one sample: after the latency, the output
        gives the integer coefficients multiplied by 8191 (impulse response);
  \item three sines of amplitude 4000 (1000 samples each, \SI{64}{\micro\second}):
        \SI{0.5}{\mega\hertz} (stop band), \SI{2.25}{\mega\hertz} (pass band)
        and \SI{5}{\mega\hertz} (stop band).
\end{itemize}
The simulation lasts \SI{210}{\micro\second} (\texttt{run 210 us} in the Tcl
console).

\fig{sim_bandpass.png}{Simulation of the band-pass filter: impulse, then sines at 0.5, 2.25 and 5 MHz}{fig:simbp}

Figure~\ref{fig:simbp} shows the expected behaviour: the output only appears
after the 22 cycles of latency, the sine at \SI{2.25}{\mega\hertz} passes with
a large amplitude, and the sines at \SI{0.5}{\mega\hertz} and
\SI{5}{\mega\hertz} are almost completely removed. At the start and at the end
of each sine there is a short transient of 31 samples (about
\SI{2}{\micro\second}), the time needed to fill the delay line. The strange
look of the \SI{5}{\mega\hertz} input comes from the display: there are only
about 3 samples per period.

\begin{table}[H]
\centering
\begin{tabular}{cccc}
\toprule
Frequency & Band & Gain (integer coefficients) & Output amplitude for $|x| = 4000$ \\
\midrule
\SI{0.5}{\mega\hertz}  & stop & $\approx 3.4$   & $\approx 14\,000$ \\
\SI{2.25}{\mega\hertz} & pass & $\approx 512$   & $\approx 2\,000\,000$ \\
\SI{5}{\mega\hertz}    & stop & $\approx 0.1$   & $\approx 400$ \\
\bottomrule
\end{tabular}
\caption{Expected gains of the band-pass filter (Octave, quantized coefficients)}\label{tab:bpgain}
\end{table}

\paragraph{Output on 14 bits.}
The coefficients are multiplied by $2^9 = 512$, so in the pass band the full
precision output is about 512 times the input (\num{4000} gives about
\num{2000000}). To get a 14-bit output we must divide by $2^9$, so we keep the
bits \texttt{data\_out(22 downto 9)}. The 9 LSBs are only the fractional part,
and bit 23 is a copy of the sign bit for normal signals. With this choice the
gain in the pass band is about 1 (0~dB), as in Octave. If we kept the 14 MSBs
(\texttt{data\_out(23 downto 10)}), we would divide by $2^{10}$ and the gain
would drop to about $1/2$ ($-6$~dB).

\paragraph{Comparison with the low-pass filter.}
The principle is the same as for the low-pass filter: coefficients on 8 bits
scaled by 512, and output taken from bit 22 down to bit 9
(\texttt{sum(22 downto 9)}). The difference is the worst case. For the
low-pass filter all the coefficients are positive and their sum is exactly
512: the gain is at most 1 and the output can never overflow 14 bits. For the
band-pass filter the coefficients have both signs: the sum of their absolute
values is 659, so the worst-case gain is $659/512 \approx 1.29$. For a sine
the gain stays at or below 1 and 14 bits are enough, but a specially built
signal could saturate slightly. Finally, with the IP, Vivado chooses the
fixed-point format and the output width (24 bits) itself, whereas for the
low-pass filter we did this work by hand.

% =====================================================================
\section{PS/PL interaction: switching the two filters}

\subsection{Block design}
\todo{TP\_M2\_source project, the two filters, the switching .vhd file, slices,
clk\_wiz at $f_s$; bus mapping: 16 LSBs = channel A, 16 MSBs = channel B, sign
extension on the 2 MSBs}

\fig{block_design_fir.png}{Block design with the two filters}{fig:bdfir}

\subsection{Switching file}
\todo{GPIO bit 0 = '0': low-pass on DAC A, band-pass on DAC B; '1': swapped}

\subsection{Programming and devmem}
\begin{lstlisting}[style=sh, caption={Loading the bitstream and switching the filters}]
bootgen -image file.bif -arch zynq -process_bitstream bin -w on
# on the Red Pitaya:
echo design_1_wrapper.bit.bin > /sys/class/fpga_manager/fpga0/firmware
devmem 0x41200004 w 0     # GPIO as output
devmem 0x41200000 w 0     # low-pass on OUT1, band-pass on OUT2
devmem 0x41200000 w 1     # swapped
\end{lstlisting}

\subsection{Tests with the generator}
\todo{cutoff frequencies measured on the board for both filters; look of the
curves at high frequency}

\fig{board_bandpass.png}{Band-pass filter on the board}{fig:boardbp}

% =====================================================================
\section{Conclusion}
\todo{summary and what we learned}

\section*{Attached files}
\begin{itemize}
  \item \texttt{filter.m}: low-pass coefficients
  \item \texttt{fir.vhd}, \texttt{tb\_fir.vhd}: low-pass filter and testbench
  \item encapsulating file and \texttt{ports.xdc}
  \item band-pass coefficients script and Vivado testbench
  \item switching \texttt{.vhd} file
\end{itemize}

\end{document}
